\unitopen{25}
\begin{candos}
\item uzasadnić działanie celem (\de{damit}, \de{um \dots{} zu}) i wskazać sprzeczność (\de{obwohl});
\item ułożyć zdania czasowe: \de{als}, \de{wenn}, \de{bevor}, \de{nachdem}, \de{während}, \de{bis}, \de{sobald};
\item odróżnić \de{als} od \de{wenn} oraz \de{damit} od \de{um \dots{} zu};
\item opowiedzieć o kolejności zdarzeń.
\end{candos}

\section{Słownictwo}
\vocabtable{25}

\section{Spójniki czasu, celu i przeciwstawienia}
Wszystkie wymagają \emph{czasownika na końcu}, jak \de{weil} z jednostki 18.

\begin{center}\small
\begin{tabular}{@{}lll@{}}
\toprule
\textsf{\textbf{Spójnik}} & \textsf{\textbf{Funkcja}} & \textsf{\textbf{Przykład}}\\
\midrule
\de{obwohl} & przeciwstawienie & \de{Obwohl ich müde \textbf{bin}, lerne ich Deutsch.}\\
\de{damit} & cel, różne podmioty & \de{Ich erkläre es, damit du es \textbf{verstehst}.}\\
\de{um \dots{} zu} & cel, ten sam podmiot & \de{Ich lerne Deutsch, um in der Schweiz zu \textbf{arbeiten}.}\\
\de{während} & jednoczesność & \de{Während ich \textbf{koche}, hört er Musik.}\\
\de{bevor} & przed & \de{Bevor ich \textbf{gehe}, schalte ich den Server aus.}\\
\de{nachdem} & po & \de{Nachdem ich \textbf{gegessen habe}, gehe ich spazieren.}\\
\de{bis} & do momentu & \de{Ich warte, bis du \textbf{kommst}.}\\
\de{sobald} & jak tylko & \de{Sobald der Test \textbf{fertig ist}, melde ich mich.}\\
\bottomrule
\end{tabular}
\end{center}

\begin{zasada}[\de{als} czy \de{wenn}?]
\de{als} = \emph{jednorazowe} zdarzenie w przeszłości (\de{Als ich in Zürich ankam, regnete es.}). \de{wenn} = zdarzenie
\emph{powtarzalne} w przeszłości lub \emph{teraźniejszość/przyszłość} (\de{Wenn ich nach Zürich fahre, nehme ich den Zug.}
\de{Wenn es regnete, blieb ich zu Hause.}). \textbf{Nachdem} wymaga czasu wcześniejszego: Perfekt w zdaniu podrzędnym przy Präsens
w głównym (\de{Nachdem ich gegessen habe, gehe ich}) lub Plusquamperfekt przy Präteritum.
\end{zasada}

\begin{falle}
Polskie „żeby” obejmuje dwie funkcje: z tym samym podmiotem używasz \de{um \dots{} zu} (bez \de{damit}):
\de{Ich lerne, \textbf{um} zu arbeiten} (nie \de{damit ich arbeite}). Różne podmioty: \de{damit}. Błąd: \de{Ich lerne, damit ich
in der Schweiz arbeiten zu können}. Drugi: \de{obwohl} i \de{trotzdem} nie są zamienne: \de{obwohl} wprowadza zdanie podrzędne,
\de{trotzdem} jest przysłówkiem (po nim odwrócony szyk: \de{Ich bin müde, trotzdem lerne ich}).
\end{falle}

\begin{szwajcaria}
Szwajcarzy często używają \de{sobald} w zapowiedziach: \de{Sobald ich Neuigkeiten habe, melde ich mich.} Standardowa formuła
końcowa w mailu: \de{Ich melde mich, sobald ich mehr weiß.} Odpowiedź w ciągu jednego dnia roboczego jest uważana za
profesjonalną, a spóźnienie bez uprzedzenia za nieuprzejme.
\end{szwajcaria}

\begin{praca}
\de{Bevor wir deployen, laufen alle Tests.} \de{Ich schreibe die Doku, damit das Team den Prozess versteht.}
\de{Um die Latenz zu senken, haben wir einen Cache eingeführt.} \de{Obwohl der Build grün war, gab es einen Fehler in
Produktion.} \de{Nachdem wir den Fehler behoben hatten, haben wir ein Postmortem gemacht.}
\end{praca}

\section{Dialogi}
\textbf{Dialog 1 --- plan wdrożenia}
\begin{dialog}
\turn{A}{Wann können wir deployen?}{Kiedy możemy wdrażać?}
\turn{B}{Sobald die Tests durchgelaufen sind. Bevor wir starten, muss das Backup fertig sein.}{Jak tylko testy przejdą. Zanim zaczniemy, backup musi być gotowy.}
\turn{A}{Und warum machen wir das heute, obwohl es Freitag ist?}{A dlaczego robimy to dziś, chociaż jest piątek?}
\turn{B}{Damit der Kunde am Montag die neue Version nutzen kann.}{Żeby klient mógł w poniedziałek używać nowej wersji.}
\end{dialog}

\section{Ćwiczenia}
\begin{cwiczenia}
\exitem{Połącz (\de{obwohl}): \de{Ich bin müde. Ich lerne.}}{\de{Obwohl ich müde bin, lerne ich.}}
\exitem{Połącz (\de{damit}): \de{Ich erkläre es. Du verstehst es.}}{\de{Ich erkläre es, damit du es verstehst.}}
\exitem{Połącz (\de{um \dots{} zu}): \de{Ich lerne Deutsch. Ich arbeite in der Schweiz.}}{\de{Ich lerne Deutsch, um in der Schweiz zu arbeiten.}}
\exitem{Połącz (\de{bevor}): \de{Ich gehe. Ich schalte den Server aus.}}{\de{Bevor ich gehe, schalte ich den Server aus.}}
\exitem{Połącz (\de{nachdem}, Perfekt): \de{Ich esse. Ich gehe spazieren.}}{\de{Nachdem ich gegessen habe, gehe ich spazieren.}}
\exitem{\de{als} czy \de{wenn}: \dots{} \de{ich in Zürich ankam, regnete es.}}{\de{Als}}
\exitem{\de{als} czy \de{wenn}: \dots{} \de{ich nach Zürich fahre, nehme ich den Zug.} (zawsze)}{\de{Wenn}}
\exitem{Popraw: \de{Ich lerne, damit ich in der Schweiz arbeiten.}}{\de{Ich lerne, um in der Schweiz zu arbeiten.}}
\exitem{Popraw: \de{Trotzdem ich müde bin, lerne ich.}}{\de{Obwohl ich müde bin, lerne ich.} lub \de{Ich bin müde, trotzdem lerne ich.}}
\exitem{Przetłumacz: „Poczekam, aż skończysz.”}{\de{Ich warte, bis du fertig bist.}}
\exitem{Przetłumacz: „Podczas gdy gotuję, słucham muzyki.”}{\de{Während ich koche, höre ich Musik.}}
\exitem{Przetłumacz: „Powiadomię Cię, jak tylko będę wiedział.”}{\de{Ich melde mich, sobald ich es weiß.}}
\end{cwiczenia}

\section{Trening słownictwa (generowany)}
\drill{25}

\weekend{Napisz pięć zdań o swoim celu zawodowym z użyciem \de{um \dots{} zu}, \de{damit}, \de{obwohl}, \de{bevor}, \de{nachdem}.
Nagraj 90 sekund: opowiedz kolejność kroków w przeprowadzce lub wdrożeniu, z czasem.}
